\documentclass[12pt]{article}
\usepackage[margin=1in]{geometry}
\usepackage[all]{xy}


\usepackage{amsmath,amsthm,amssymb,color,latexsym}
\usepackage{geometry}        
\geometry{letterpaper}  
\usepackage{boxedminipage}
\usepackage{graphicx}
\usepackage{hyperref}
\hypersetup{
    colorlinks=true,
    citecolor=blue,
    linkcolor=blue,
    filecolor=cyan,
    urlcolor=magenta,
}
\usepackage[capitalize,nameinlink]{cleveref}

\newtheorem{problem}{Problem}

\newenvironment{solution}[1][\it{Solution}]{\textbf{#1. } }{$\square$}
\newcommand{\pointsbox}[1]{%
  \colorbox{red!70}{\strut\textcolor{white}{\sffamily\bfseries\ \ #1\ \ }}%
}
\newcommand{\pagetag}[1]{%
  \colorbox{red!70}{\strut\textcolor{white}{\sffamily\bfseries\ \ #1\ \ }}%
}
\newcommand{\striptitle}[2]{%
  \noindent
  \begin{minipage}{0.07\linewidth}\centering \pagetag{#1}\end{minipage}%
  \hfill
  \begin{minipage}{0.89\linewidth}
    {\large\scshape #2}
  \end{minipage}\par\vspace{0.75em}
}
% --- Environments ---
\theoremstyle{plain}
\newtheorem{theorem}{Theorem}
\theoremstyle{definition}
\newtheorem{definition}[theorem]{Definition}
\newtheorem{remark}[theorem]{Remark}

\newenvironment{boxedalgo}
  {\begin{center}\begin{boxedminipage}{1\linewidth}}
  {\end{boxedminipage}\end{center}}

\newenvironment{inlinealgo}[1]
  {\noindent\hrulefill\\\noindent{\bf {#1}}}
  {\hrulefill}


% --- Shortcuts ---
\newcommand{\R}{\mathbb{R}}
\newcommand{\inner}[2]{\left\langle #1,#2 \right\rangle}
\newcommand{\proj}{\mathrm{proj}}
\newcommand{\Id}{\mathbf{I}}
\newcommand{\ot}{\otimes}
\newcommand{\mc}{\mathcal}

\begin{document}
\noindent CS 578 Statistical Machine Learning (Fall 2026)\hfill Problem Set \# 1\\
%\rightline{\textbf{Due: September 15th}}
\{\{Your name\}\} \hfill \textbf{Due: September 6th}

\noindent\hrulefill

\subsection*{Problem 1. Bayes predictor.}
The Bayes predictor $f_{\star}:{\cal X}\rightarrow {\cal Y}$ is defined as:
\begin{align*}
    f_{\star}(x')\in \arg \min_{z\in {\cal Y}} \mathbb{E}[\ell(y,z)|x=x'].
\end{align*}
\begin{enumerate}
    \item Consider binary classification with ${\cal Y} = \{-1, 1\}$ with the loss function 
        \begin{align*}
            \ell(y,z)=\begin{cases}
            0 & \text{if}~y=z,\\
            c_{\rm FP}>0 & \text{if}~y=-1,z=1,\\
            c_{\rm FN} > 0 & \text{if}~y=1,z=-1.
            \end{cases}
        \end{align*}
    Compute a Bayes predictor at $x=x'$ as a function of $\mathbb{E}[y|x=x']$.\\
    \textit{Hint: Show that $f_\star(x') = {\bf 1}\{\mathbb{E}[y|x=x'] \geq~\text{or }\leq~\text{???}\}$}
    \item Show that every Bayes predictor minimizes the population risk; that is, $\mathcal{R}(f_\star)\leq \mathcal{R}(g)$ for any $g:{\cal X}\rightarrow {\cal Y}$.
    \item \emph{(Bonus)} We consider a learning problem on ${\cal X} \times {\cal Y}$, with ${\cal Y} = \mathbb{R}$ and the absolute loss defined as $\ell(y, z) = |y - z|$. Characterize the Bayes predictors in terms of the conditional distribution of $y|x=x'$, and give one explicit choice using the conditional cumulative distribution function (CDF).
\end{enumerate}
\noindent{\bfseries Solution.}\newline 
\   %%% <-- ERASE THIS LINE AND WRITE YOUR SOLUTION HERE
\newpage

\subsection*{Problem 2. Empirical risk minimization and PAC learning.} 
Let ${\cal X}$ be a discrete (finite or infinite) domain, ${\cal Y}=\{0,1\}$, and we use the 0--1 loss. The training sample consists of $m$ i.i.d. samples drawn from a distribution \emph{realizable} by ${\cal H} = \{h_z:z \in {\cal X}\} \cup \{h_-\}$, where for each $z \in {\cal X}$, $h_z$ is the indicator function for $z$:
\begin{align*}
    h_z(x)=\begin{cases}
        1 & \text{if}~x=z,\\
        0 & \text{otherwise}.
    \end{cases}
\end{align*}
$h_-$ is the all-zero function, namely, $\forall x \in  {\cal X}$, $h_-(x) = 0$. 
\begin{enumerate}
    \item Describe an algorithm that implements the ERM rule for learning ${\cal H}$.
    \item Show that ${\cal H}$ is PAC learnable. Provide an upper bound on the sample complexity.
\end{enumerate}
\noindent{\bfseries Solution.}\newline 
\   %%% <-- ERASE THIS LINE AND WRITE YOUR SOLUTION HERE
\newpage

\subsection*{Problem 3. Probability and concentration.} 
Suppose that we have $m$ coins. Each coin has the same probability of landing on heads, denoted $p$. Suppose we pick up each of the $m$ coins in turn and, for each coin, perform $n$ tosses. We assume that those $mn$ tosses are mutually independent. Note that the probability of obtaining exactly $k$ heads out of $n$ tosses for any given coin is given by the binomial distribution: 
\begin{align*}
    \binom{n}{k} p^{k} (1-p)^{n-k}.
\end{align*}
Define
\begin{align*}
    \hat{p}_i:=\frac{\text{number of times coin}~i~\text{lands on heads}}{n}.
\end{align*}
\begin{enumerate}
    \item Compute a formula for the exact probability that at least one coin will have $\hat{p}_i=0$.
    \item For $\epsilon >0$,  apply Hoeffding’s inequality\footnote{\url{https://en.wikipedia.org/wiki/Hoeffding\%27s_inequality}} and a union bound to prove an upper bound on
    \begin{align*}
        \Pr\left[\max_{1\leq i\leq m} |\hat{p}_i-p|>\epsilon\right]
    \end{align*}
\end{enumerate}
\noindent{\bfseries Solution.}\newline 
\   %%% <-- ERASE THIS LINE AND WRITE YOUR SOLUTION HERE
\newpage



\end{document}
